\chapter{Implementation Details and Simulation Code}
\label{appendix:code}

This appendix provides the Python implementation of the main components of this thesis. The code is organized into three main sections: time and frequency domain validation (Chapter 5), the map-making and harmonic analysis pipeline (Chapter 6), and functions optimized using Just-In-Time (JIT).

\section{Time and Frequency Domain Analysis}
This script handles noise model validation by producing simulated time-ordered data (TOD) and the theoretical power spectral density (PSD) model. Although the following code uses a "mock detector" for demonstration, a real detector can be loaded from the LiteBIRD instrument model (IMo).

\begin{lstlisting}[language=Python, caption={Validation of noise statistics in time and frequency domains.}]
import litebird_sim as lbs
import numpy as np
from scipy.signal import welch
from mylib import add_lbsNoise, add_OofaNoise

# --- Simulation Setup ---
start_time = 0
time_span_s = 1000.0  # Duration sufficient for 1/f characterization

sim = lbs.Simulation(
    start_time=start_time,
    duration_s=time_span_s,
    random_seed=12345,
    imo=lbs.Imo(flatfile_location=lbs.PTEP_IMO_LOCATION)
)

# --- Scanning Strategy ---
sim.set_scanning_strategy(
    lbs.SpinningScanningStrategy(
        spin_sun_angle_rad=np.deg2rad(0),
        precession_rate_hz=0,
        spin_rate_hz=1/60,
        start_time=start_time,
    ),
    delta_time_s=5.0,
)

# --- Detector Setup ---
# Note: Parameters can be substituted with IMO data 
det = lbs.DetectorInfo(
    name="Boresight_detector",
    sampling_rate_hz=20.0,
    bandcenter_ghz=100.0,
    net_ukrts=50.0,
    fknee_mhz=500.0,
    fmin_hz=1e-5,
    alpha=1.0
)

sim.create_observations(detectors=det)
tod = sim.observations[0].tod 

# --- Noise Generation ---
# Model selection: add_lbsNoise (FFT-based) or add_OofaNoise (Filter-based)
add_lbsNoise(tod, det, block_duration_s=100)

# --- Spectral Estimation (Welch Method) ---
frequencies, psd = welch(
    tod,
    fs=det.sampling_rate_hz,
    window='hann',
    nperseg=2**12,
    noverlap=2**11,
    scaling='density',
    return_onesided=False
)

# Sorting frequencies for analysis
sort_idx = np.argsort(frequencies)
frequencies = frequencies[sort_idx]
psd = psd[:, sort_idx]

# --- Theoretical PSD Comparison ---
f = frequencies
f_knee = det.fknee_mhz / 1000.0  # Convert to Hz
f_min = det.fmin_hz

# Analytical model for 1/f + white noise
theoretical_psd = (det.net_ukrts**2) * ( (f**det.alpha + f_knee**det.alpha) / (f**det.alpha + f_min**det.alpha) )
\end{lstlisting}





\section{Full Map-Making and Harmonic Analysis Pipeline}
This section presents the integrated pipeline used to generate the long-term simulation results discussed in Chapter 6. The implementation adopts an incremental map-making strategy to maintain a minimal memory footprint by processing 365 days of observation in 12-hour chunks. 

The final part of the script illustrates the harmonic decomposition of the resulting noise maps to estimate the angular power spectrum $C_\ell^{TT}$.

\begin{lstlisting}[language=Python, caption={Incremental map-making and angular power spectrum estimation.}]
import healpy as hp
import numpy as np
from uuid import UUID
import litebird_sim as lbs
from mylib import add_OofaNoise, add_lbsNoise, add_samples_to_map, normalize_map

# --- 1. Simulation & Instrument Configuration ---
nside = 128  # Map resolution
lmax = 3 * nside - 1
chunk_length = 3600 * 12  # 12-hour blocks
number_of_days = 365
duration_s = 86400 * number_of_days 
num_of_observations = duration_s // chunk_length

sim = lbs.Simulation(
    start_time=0,
    duration_s=duration_s,
    random_seed=12345,
    imo=lbs.Imo(flatfile_location=lbs.PTEP_IMO_LOCATION)
)

# Scanning Strategy (LiteBIRD Spinning Strategy)
sim.set_scanning_strategy(
    lbs.SpinningScanningStrategy.from_imo(
        imo=lbs.Imo(lbs.PTEP_IMO_LOCATION), 
        url=UUID("117fd641-a925-4eeb-9fda-f7c468682108")
    ),
    delta_time_s=120.0
)

# Instrument and Detector Setup
instrument = lbs.InstrumentInfo(
    boresight_rotangle_rad=0.0,
    spin_boresight_angle_rad=np.deg2rad(90),
    spin_rotangle_rad=np.deg2rad(75)
)
sim.set_instrument(instrument)

det = lbs.DetectorInfo(
    name="Boresight_detector", sampling_rate_hz=20.0, net_ukrts=50.0,
    fknee_mhz=500.0, fmin_hz=1e-5, alpha=1.0
)

# --- 2. Incremental Processing Loop ---
# Allocate metadata without full TOD arrays to save RAM
sim.create_observations(detectors=det, num_of_obs_per_detector=num_of_observations, allocate_tod=False)

# Pre-allocate map buffers and temporary TOD block
accum_map = np.zeros(hp.nside2npix(nside))
hit_count_map = np.zeros(len(accum_map), dtype='int64')
tod_buffer = np.empty((1, int(det.sampling_rate_hz * chunk_length)))

for obs_idx, cur_obs in enumerate(sim.observations):
    # A. Pointing Generation
    lbs.prepare_pointings(cur_obs, instrument, sim.spin2ecliptic_quats)
    pointings, _ = cur_obs.get_pointings("all")
    pixidx = hp.ang2pix(nside, pointings[0,:,0], pointings[0,:,1])
    
    # B. Noise Generation (In-place)
    tod_buffer[:, :] = 0.0
    add_OofaNoise(tod_buffer, det, block_duration_s=chunk_length)
    
    # C. Map Accumulation (using Numba-optimized kernels)
    add_samples_to_map(tod_buffer, pixidx, accum_map, hit_count_map)  

# Final normalization (Binned Map-Making)
normalize_map(accum_map, hit_count_map)

# --- 3. Harmonic Analysis (Power Spectrum Estimation) ---
# The following steps decompose the map into Spherical Harmonics 
# and compute the auto-correlation power spectrum.

almT = hp.map2alm(accum_map, lmax=lmax)
clTT = hp.alm2cl(almT)
\end{lstlisting}





\section{Numba-Optimized Kernels}
These functions are decorated with \texttt{@njit} to handle the high-speed accumulation of TOD samples into HEALPix pixels.

\begin{lstlisting}[language=Python, caption={Performance-critical Numba kernels for sample accumulation and map normalization.}]
from numba import njit

@njit
def add_samples_to_map(tod, pixel_indexes, accum_map, hit_count_map):
    """
    Accumulates TOD samples into the HEALPix maps using JIT-optimized loops.
    
    Parameters:
    - tod: 2D array (ndet, nsample) containing the noise data.
    - pixel_indexes: 1D array of pixel indices corresponding to the scan strategy.
    - accum_map: The map where signal values are summed.
    - hit_count_map: The map counting the number of observations per pixel.
    """
    ndet, nsample = tod.shape

    for det_idx in range(ndet):
        for sample_idx in range(nsample):
            pix_idx = pixel_indexes[sample_idx]
            accum_map[pix_idx] += tod[det_idx, sample_idx]
            hit_count_map[pix_idx] += 1

@njit
def normalize_map(accum_map, hit_count_map):
    """
    Performs binned map-making normalization by dividing the accumulated 
    signal by the hit count.
    """
    for pix_idx in range(len(accum_map)):
        if hit_count_map[pix_idx] > 0:
            accum_map[pix_idx] /= hit_count_map[pix_idx]
\end{lstlisting}
